\documentclass[11pt,a4paper]{ctexart}
\usepackage[margin=18mm,headheight=15pt]{geometry}
\usepackage{amsmath,amssymb,mathtools,bm,longtable,array,enumitem,xcolor,hyperref,fancyhdr,url}
\hypersetup{colorlinks=true,linkcolor=blue!55!black,pdfauthor={CHASI research notes},pdftitle={茶思研究汇报现场答辩速查}}
\setlength{\parindent}{2em}
\setlength{\parskip}{4pt}
\setlength{\emergencystretch}{3em}
\allowdisplaybreaks
\pagestyle{fancy}\fancyhf{}\fancyfoot[C]{\thepage}
\begin{document}
\begin{center}\LARGE\bfseries 茶思研究汇报现场答辩速查\end{center}

\begin{quote}\small 对应 15 页用户当前 PPTX（SHA-256 前缀 \path{f460c1bcdd725b24}）；正式代码提交 \path{db9d7e3}。完整校验值及逐式推导见《CHASI\_逐页汇报与答辩讲解稿》。\end{quote}

\subsection*{30 秒主线}

原始茶思任务要求在七站三扇区、非理想 SRS/CSI、273 RB、每 TRP 功率约束下研究用户中心多 TRP 协同的边缘体验与规模。当前 MATLAB 场景生成 $H^{\rm true}$；等效观测产生 $\widehat H,C$；规则固定服务、调度、流数和资源权重；集中式矩阵 WMMSE–ADMM 求完整波束 $w$；同一 $w$ 在真信道上评价。202610 方案拟研究校准的非均匀误差椭球、最坏情形、稀疏/休眠与事件触发；这些不是当前已运行结果。\par

\subsection*{沿 PPT 页码的一句话}

\begin{itemize}[leftmargin=2em,itemsep=1pt]
\item 1–3：问题是测量不准、边缘干扰、跨组功率与计算规模；右侧指标是检验量，不是结果。
\item 4：7 站 × 3 共址 TRP，当前每小区 10 UE；彩色点是真实 UE，颜色是初始归属站址；UE7/T01/T02 是第 1 组示例。
\item 5–6：$H^{\rm true}\to(\widehat H,C)\to(b,\chi,r,\nu)\to w\to R^{\rm true}$；测量/CPU 泳道是等效仿真，不是完整协议与分布式实现。
\item 7–8：历史 GA/PSO 链未接正式入口；当前只优化 $w$；真信道速率再进入静态业务体验\textbf{代理}。
\item 9–10：外轮重算 $F,R^{yy},A,E,\Omega$；内轮 CG 解 $x$、联合功率投影 $z$、更新缩放对偶 $y$；返回可行 $z$。
\item 11–12：$C$ 进 $R^{yy}$ 和二次算子；$\Omega$ 是流内 MSE 权重，$\eta$ 是用户间效用权重。
\item 13–15：原文 $64$ 维单接收向量椭球与当前 $4\times64$ 信道/$4\times4$ 接收侧统计不同；最坏情况、稀疏、休眠、TRP 本地计算、触发均属拟研究。
\end{itemize}

\subsection*{必背公式与维度}

\[
H^{\rm true}_{klg}=\sqrt{\beta_{kl}}L_{\rm rx}Z_{klg}\in\mathbb C^{4\times64},
\quad L_{\rm rx}\in\mathbb C^{4\times4},
\quad Z_{klg}\in\mathbb C^{4\times64}.
\]
\[
F_{kjg}=\sum_l\widehat H_{klg}w_{ljg}\in\mathbb C^{4\times r_{jg}},
\quad p_{lg}=\sum_j\|w_{ljg}\|_F^2,
\quad R^{yy}_{kg}=I_4+\sum_jF_{kjg}F_{kjg}^H+\sum_lp_{lg}C_{klg}.
\]
\[
A_{kg}=(R^{yy}_{kg})^{-1}F_{kkg}\in\mathbb C^{4\times r},
\quad E_{kg}=I_r-A^HF_{kkg}-F_{kkg}^HA+A^HR^{yy}A,
\quad\Omega_{kg}=E_{kg}^{-1}\in\mathbb C^{r\times r}.
\]
\[
\sum_g\nu_g\sum_j\|w_{ljg}\|_F^2\le P_l,
\quad(\nu_g\mathcal Q_g+\tfrac\rho2 I)x_g=\nu_g B_g+\tfrac\rho2(z_g-y_g),
\quad z=\Pi_{\mathcal P}(x+y),\quad y^+=y+x-z.
\]
\[
R_{kg}^{\rm SE,true}=\log_2\det(I_4+\sum_jF^{\rm true}_{kjg}(F^{\rm true}_{kjg})^H)
-\log_2\det(I_4+\sum_{j\ne k}F^{\rm true}_{kjg}(F^{\rm true}_{kjg})^H).
\]

\subsection*{五个最容易被追问的边界}

1. \textbf{$R^{yy}$ 的 $j$ 包含 $k$；干扰分母的 $j\ne k$ 排除 $k$。} 同一用户跨 TRP 先在 $F$ 中相干加幅度。\par
2. \textbf{默认 $I_4$ 是归一化噪声。} $w$ 含发射幅度；不能再无定义地乘 $\sqrt p$。\par
3. \textbf{功率可行 $\ne$ ADMM 收敛 $\ne$ WMMSE 停止。} 保存对拍也只证明同输入数值一致；正式规模停止条件未满足。\par
4. \textbf{默认静态评价没有队列出队。} 页面“队列限制”“按协议体验率”宜口头改成忙闲/尾段门控的容量与体验代理；原图协议需再核对 burst 与 bottom 5\% 样本分母。\par
5. \textbf{$C\in\mathbb C^{4\times4}$ 不是原文 $Q\in\mathbb C^{64\times64}$。} 原文最坏情况式只有研究目标，尚无矩阵信道下的校准、转化与效果证据。\par

\subsection*{研究计划回答模板}

\textbf{已有基础：} 矩阵 WMMSE、接收侧链路条件误差统计与集中式 ADMM。\par
\textbf{切入点：} 误差尺度/方向及链路服务与整 TRP 激活的数学对象需明确。\par
\textbf{拟研究：} 校准全矩阵误差集合、推导最坏情形内层；定义链路组稀疏和休眠后可行重优化；拆分跨 TRP 耦合并定义消息缓存/触发量；研究慢层候选与快层统一评价接口。\par
\textbf{验证：} 固定场景、业务、预算、容差和随机流，做名义/统一/非均匀对照；分别报告边缘体验、空口功率、可行性、内外残差、时间与消息量，不预设增益。\par
\end{document}
